\section{AutoEncoder Modulation as a Supporting Study}
AutoEncoder-based Modulation is included in the thesis because the LDPC decoder operates on the reliability information produced by the receiver front end. If the Modulation and Demodulation process produces biased LLRs, the LDPC decoder may not reach its expected performance even when the decoder algorithm is well designed. Therefore, a study of learned Modulation is relevant as a supporting analysis of the soft-information source in a BIBCM-ID-oriented receiver. The key point is that AutoEncoder Modulation is not presented as the main LDPC-decoding contribution. It is used to examine how a trainable signal mapping and a trainable receiver can adapt to channel impairments that weaken conventional Demodulation.

In a communication AutoEncoder, the transmitter network maps a message or bit pattern to a channel input:
\begin{equation}
  \mathbf{x}=f_\theta(\mathbf{m}),
\end{equation}
the channel produces
\begin{equation}
  \mathbf{y}=\mathcal{H}(\mathbf{x})+\mathbf{w},
\end{equation}
and the receiver network estimates either the transmitted message or the transmitted bits:
\begin{equation}
  \hat{\mathbf{p}}=g_\phi(\mathbf{y}).
\end{equation}
The parameters $\theta$ and $\phi$ are learned by minimizing a loss such as cross-entropy. This formulation resembles an end-to-end communication system, but the interpretation must be precise. If the receiver outputs only a message class, the model is a learned detector. If the receiver outputs bit probabilities or LLR-like values, it becomes more directly relevant to an iterative coded-modulation receiver.

Power normalization is central to a fair comparison. Without a constraint, the transmitter network could reduce the error rate by increasing the signal energy. A typical normalization is
\begin{equation}
  \mathbf{x}_{\mathrm{norm}}
  =
  \sqrt{P_0}
  \frac{\mathbf{x}_{\mathrm{raw}}}
  {\sqrt{\mathbb{E}\{\|\mathbf{x}_{\mathrm{raw}}\|^2\}}},
\end{equation}
which enforces an average-power constraint. When comparing AutoEncoder Modulation with 16-QAM or another baseline, the same effective energy normalization must be used. Otherwise, a BER improvement may reflect higher transmitted power rather than a better signal representation.

The channel layer determines what the AutoEncoder learns. Under AWGN only, a learned constellation may resemble a conventional constellation because Euclidean distance is the dominant design criterion. Under PA nonlinearity, the learned constellation may avoid high-amplitude regions that suffer compression. Under CFO, a sequence-based transmitter and receiver may learn a representation that is robust to phase rotation over time. Under fading and ISI, the receiver may learn equalization-like behavior. These observations explain why the AutoEncoder results are useful for discussing Modulation and Demodulation, even though they are not the central evidence for ANN-assisted LDPC decoding.

\section{PA Nonlinearity and Learned Signal Geometry}
Power-amplifier nonlinearity is one of the most important reasons for studying learned Modulation. Practical transmitters often operate power amplifiers near saturation to improve energy efficiency. Near saturation, the amplifier no longer behaves as a linear gain. A common memoryless model expresses the output as
\begin{equation}
  z=A(r)e^{j(\angle x+\Phi(r))}, \qquad r=|x|,
\end{equation}
where $A(r)$ is the AM/AM conversion and $\Phi(r)$ is the AM/PM conversion. For a Rapp-type AM/AM model,
\begin{equation}
  A(r)=\frac{r}{\left(1+\left(\frac{r}{A_{\mathrm{sat}}}\right)^{2p}\right)^{1/(2p)}}.
\end{equation}
When $r$ is small, $A(r)\approx r$. When $r$ approaches the saturation region, $A(r)$ grows slowly and high-amplitude points are compressed. For square QAM, outer constellation points are affected more strongly than inner points. The received clusters are no longer located at the points assumed by a linear-channel Demodulator.

The effect on LLRs can be severe. A conventional Demodulator may compute
\begin{equation}
  p(y|x)\propto \exp\left(-\frac{|y-hx|^2}{N_0}\right),
\end{equation}
while the actual received sample is closer to
\begin{equation}
  y=hA(|x|)e^{j(\angle x+\Phi(|x|))}+w.
\end{equation}
The mismatch changes both distances and decision boundaries. A learned Modulation scheme can respond by reshaping the constellation before the amplifier so that the post-amplifier points are more separable. Alternatively, a neural receiver can learn the nonlinear inverse or a corrected decision rule.

This interpretation should be kept moderate. Learned Modulation does not automatically solve PA distortion. It depends on the accuracy of the PA model used in training and on whether the operating amplifier matches that model. If the network is trained for one saturation level and tested at another, performance may degrade. Therefore, robust training should randomize the PA model parameters:
\begin{equation}
  A_{\mathrm{sat}}\sim\mathcal{U}(A_{\min},A_{\max}), \qquad
  p\sim\mathcal{U}(p_{\min},p_{\max}).
\end{equation}
This training strategy would encourage the transmitter and receiver to learn a representation that is not overly specialized to one amplifier condition. In the present work, it is treated as a design implication for future development rather than as a claim that the current simulations already cover all hardware variation.

\section{Carrier-Frequency Offset and Sequence-Based Demodulation}
Carrier-frequency offset appears when transmitter and receiver oscillators are not perfectly aligned. In discrete time, the received sample can be modeled as
\begin{equation}
  y[k]=x[k]e^{j(2\pi\epsilon k+\varphi_0)}+w[k],
\end{equation}
where $\epsilon$ is the normalized CFO and $\varphi_0$ is the initial phase. A memoryless Demodulator observes only one rotated symbol at a time. If the phase is unknown, the point may appear closer to the wrong decision region. A sequence-based receiver observes multiple samples:
\begin{equation}
  \mathbf{y}_{k:k+T-1}=[y[k],y[k+1],\ldots,y[k+T-1]],
\end{equation}
and can exploit the fact that the phase rotation evolves systematically with $k$.

This explains why the multi-symbol AutoEncoder scenario is relevant. A learned receiver with sequence input can implicitly estimate or compensate a rotation trend. The benefit is not merely that a neural network is used; the benefit comes from providing temporal context to the receiver. In classical communication systems, CFO is normally handled by synchronization algorithms. The AutoEncoder result can be interpreted as evidence that a neural receiver can absorb part of the synchronization burden under a controlled scenario.

The complexity trade-off must also be stated. Sequence processing increases latency and memory because the receiver uses a window of samples rather than a single symbol. If the sequence length is $T$ and the hidden dimension is $d_h$, the computational cost can grow approximately with $T d_h$ for simple feed-forward processing and more for recurrent or attention-based models. Therefore, a sequence-based neural receiver is useful only when the performance gain justifies the increased latency and complexity.

For BIBCM-ID relevance, CFO affects the LLRs delivered to the LDPC decoder. If CFO is not corrected, the Demodulation block may produce low-confidence or wrong-confidence messages. A learned Demodulator that is robust to CFO could provide better soft information. However, to support iterative decoding, it should output bit-level reliability values rather than only message decisions. A possible interface is
\begin{equation}
  \hat{L}_k=g_\phi(\mathbf{y}_{k:k+T-1},L_{A,k}),
\end{equation}
where $L_{A,k}$ represents a priori information from the decoder. This equation is a design direction for learned soft Demodulation; it shows how AutoEncoder ideas can be made compatible with BIBCM-ID principles.

\section{LLR Calibration for Neural Demodulation}
The most important scientific bridge between AutoEncoder Modulation and LDPC decoding is LLR calibration. A neural receiver trained with cross-entropy may output probabilities that are useful for classification but not perfectly calibrated \cite{guo2017}. Calibration means that when a receiver outputs probability $0.9$, the event should be correct approximately 90 percent of the time over many samples with similar confidence. Poor calibration is harmful for iterative decoding because the LDPC decoder uses probability magnitudes as reliability weights.

For bit $b_k$, a neural receiver may output $\hat{p}_k=P(b_k=1|\mathbf{y})$. The corresponding LLR is
\begin{equation}
  \hat{L}_k=\log\frac{\hat{p}_k}{1-\hat{p}_k}.
\end{equation}
If $\hat{p}_k$ is overconfident, $|\hat{L}_k|$ becomes too large. If the sign is wrong, a large magnitude can mislead the LDPC decoder. A simple calibration method is temperature scaling:
\begin{equation}
  \hat{p}_k(T)=\sigma\left(\frac{z_k}{T}\right),
\end{equation}
where $z_k$ is the neural logit, $\sigma(\cdot)$ is the sigmoid function, and $T>0$ is a calibration temperature. If $T>1$, the output probabilities become less confident. If $T<1$, they become sharper. The temperature can be selected on a validation set by minimizing negative log-likelihood.

Calibration can also be evaluated with the Brier score:
\begin{equation}
  \mathrm{Brier}=\frac{1}{N}\sum_{n=1}^{N}(\hat{p}_n-b_n)^2,
\end{equation}
or with reliability diagrams that compare predicted confidence with empirical accuracy. These metrics are not substitutes for BER, but they explain whether the receiver is producing usable soft information. For an iterative coded-modulation system, a receiver with slightly worse hard-decision accuracy but better calibrated LLRs may sometimes support better decoding than an overconfident classifier.

This point is particularly important for the thesis title. The central object is LDPC decoding, and LDPC decoding depends on soft information. AutoEncoder Modulation is therefore defensible as a supporting topic only when it is linked to soft-output Demodulation and trainable receiver interfaces \cite{cammerer2020}. The thesis makes that link explicit: learned Modulation is not discussed as an isolated deep-learning demonstration, but as a way to study how channel-aware signal representation and neural receiver processing may improve the reliability information supplied to the decoder.

\section{Quality Criteria for Interpreting AutoEncoder Results}
AutoEncoder communication results should be interpreted using several criteria. The first criterion is the baseline. A learned constellation should be compared with a conventional scheme under the same power constraint, channel model, SNR definition, and receiver information. If the baseline receiver lacks channel-state information while the neural receiver implicitly learns it from training, the comparison may be unfair. Therefore, the presented figures are treated as evidence of feasibility and trend, not as a general performance claim.

The second criterion is the training distribution. If the model is trained and tested on exactly the same impairment parameters, the result measures interpolation within a narrow condition. A stronger result would test random PA parameters, multiple CFO values, fading distributions not seen during training, and different SNR ranges. This does not invalidate the presented results; it defines their scope.

The third criterion is the output interface. A message-classification AutoEncoder is useful for studying learned signal geometry. A BIBCM-ID-compatible neural Demodulator must go further and produce bit-level soft information:
\begin{equation}
  \mathbf{\hat{L}}=[\hat{L}_1,\ldots,\hat{L}_m].
\end{equation}
If iterative feedback is used, the neural Demodulator should also accept a priori information from the decoder. Without this interface, the receiver remains a learned detector rather than an iterative Demodulation component.

The fourth criterion is complexity. A neural receiver with a large number of layers may improve a BER curve but may not be attractive for a practical receiver. Complexity should be evaluated through parameter count, multiply-accumulate operations, latency, memory access, and quantization robustness. These criteria mirror the ANN-assisted LDPC discussion: a neural method is most valuable when it improves or preserves performance while maintaining a realistic implementation path.

The fifth criterion is interpretability. Classical Modulation designs can be explained through Euclidean distance, Gray labeling, PAPR, and likelihood functions. Learned Modulation should be analyzed using similar concepts where possible. For example, if the learned constellation moves points away from the PA saturation region, this can be explained through AM/AM distortion. If a sequence receiver improves under CFO, this can be explained through temporal phase structure. Such interpretation makes the AutoEncoder part more suitable for an engineering thesis.

\section{Connection Between the Two Research Directions}
The two research directions in the thesis are connected by soft-information quality. The ANN-assisted LDPC decoder acts inside the channel decoder. It modifies the way reliability messages are combined under parity constraints. The AutoEncoder Modulation study acts before the decoder. It explores how the transmitted representation and receiver observation processing affect the reliability values that enter the decoder. Both directions therefore address the same system-level concern from complementary positions.

This connection can be written compactly. Let $\mathbf{L}_{\mathrm{dem}}$ denote the LLR vector produced by the Demodulation block and let $\mathcal{D}_\theta$ denote an ANN-assisted LDPC decoder. The decoded estimate is
\begin{equation}
  \hat{\mathbf{u}}=\mathcal{D}_\theta(\mathbf{L}_{\mathrm{dem}}).
\end{equation}
The decoder performance depends on both $\mathcal{D}_\theta$ and the distribution of $\mathbf{L}_{\mathrm{dem}}$. AutoEncoder Modulation changes this distribution by changing the signal mapping or receiver processing:
\begin{equation}
  \mathbf{L}_{\mathrm{dem}}=g_\phi(\mathbf{y},\mathbf{L}_A).
\end{equation}
Thus, the AutoEncoder study is relevant because it examines the upstream conditions under which the ANN-assisted LDPC decoder operates.

This connection should be interpreted without overstatement. The present analysis does not require a claim that both neural components have been jointly optimized. It is sufficient and more scientifically defensible to show that the LDPC decoder contribution is direct, while the AutoEncoder Modulation analysis identifies a compatible supporting direction. This framing resolves the apparent mismatch between the thesis title and Chapter 3. Chapter 3 is not a detour from LDPC decoding; it studies how Modulation and Demodulation affect the soft information that LDPC decoding consumes.
